% !TeX root = ../main.tex
\section{Discussion and Limitations}\label{sec:discussion}

In this section, we discuss the limitations of \SaiLoR and of this report.
For each limitation, we state what it affects and how we mitigate it.

\subsection{Limitations}\label{sec:discussion-limitations}

% No user evaluation
We have not evaluated \SaiLoR with users.
This report describes the tool and its design, but it does not show that \SaiLoR makes reviews faster, more accurate, or easier to trace than other tools.
We derived the design from the guidelines for \SLRs, and it might not match the needs of every research group.
A user study is needed to assess the benefit (cf.~\autoref{sec:conclusion}).

% Reverse-engineered requirements
The requirements were reverse-engineered from the code by the developer of the tool.
Thus, they describe what \SaiLoR does and not what users need.
If the code has a bug, a requirement can describe the buggy behavior as intended.
Also, the developer might have overlooked behavior or described it in favor of the implementation.
We mitigate this threat with the evidence of each requirement, because tests, source lines, and commits can be checked by others.
Still, an independent review of the requirements is open.

% Tauri port
The move to Tauri is recent.
The guard against reload shortcuts is now ported to Tauri as well.
However, no automated test shows its dialog in a running application, and we have not tried it by hand on Windows and Linux.
Furthermore, \autoref{fig:screening} and \autoref{fig:merge} still show the user interface of version 1.x, because the online demo has no screening project and no git view.

% Platform tests
Our tests cover the three operating systems unevenly.
The probe of the desktop application runs on macOS, Windows, and Linux, but a WebDriver session drives it only on Linux and Windows, because the Tauri driver does not support macOS.
Most behavior is tested in the unit and integration suites, which run in a simulated browser and not in the web views of the operating systems.
Differences between WKWebView, WebView2, and WebKitGTK can therefore remain unnoticed.
The \CI probe renders a \PDF in all three web views, but the probe does not yet cover further workflows.
Likewise, we checked the right-to-left layout for Persian only in the simulated browser.

% Performance numbers
The performance numbers in this report come from benchmark tests with generated graphs and from single runs on a laptop under load.
They show that the limits work, but they do not predict the speed on the computer of a reviewer.
A real review graph can have a different structure than the generated graphs.

% Mutation testing
The mutation testing covers only the model code and two modules of the \LLM assistance.
The user interface components, the store, and the Rust backend are not mutated.
Also, we never completed a full run over the model code, only runs per file.

% Strict mode
The strict mode helps a team to follow its process, but it does not enforce the process.
A person with access to the files can change any answer.
Teams that need this protection have to rely on the access rights of their git hosting.

% Unsigned builds
The installers are not signed with a certificate from Apple or Microsoft.
macOS and Windows warn users the first time they open \SaiLoR, and users have to confirm the start.
This step can confuse users and can also train them to ignore such warnings.
The README explains the steps for each operating system.
Updates are signed with a key of the project, and \SaiLoR verifies the signature before it installs an update.

% Local LLM runtime
We verified the managed local runtime for decision models only on macOS on arm64.
The code paths for \GPUs with CUDA and Vulkan have not run on real hardware.
In addition, the managed installation stays disabled until llama.cpp releases the route that the decision models need.
Until then, users who want local decision models have to run a server themselves.
Hosted providers and local chat servers are not affected by this limitation.

% AI preview
The \LLM assistance is a preview.
We have not measured how accurate its proposals are, and studies report errors of \LLMs in screening and extraction~\cite{felizardo2024chatgpt,khraisha2024llm}.
A reviewer who accepts proposals without checking them can bring wrong data into the review.
\SaiLoR reduces this risk.
It checks the quoted evidence, marks the fields that the model wrote, and records each use for the report.
It cannot replace the judgment of the reviewer.

% Scope
\SaiLoR does not cover all activities of an \SLR.
It records the search strings and databases, but it does not run searches.
It also does not synthesize the extracted data.
The knowledge graph view and its queries help to explore the data, but reviewers still need their own scripts or tools for the analysis.

\subsection{Availability}\label{sec:discussion-availability}

% Release
\SaiLoR is open source under the GPL-3.0 license.
The source code is available on GitHub.\footnote{\url{https://github.com/Gram21/SaiLoR}}
The release workflow builds installers for macOS 15 or later on Apple silicon and Intel, for Windows, and for Linux on x86-64 and arm64 as AppImage, Debian, and RPM packages.
The Zenodo concept \DOI~\cite{sailor2026zenodo} resolves to the latest archived version.
A project site on GitHub Pages hosts the online demo.\footnote{\url{https://gram21.github.io/SaiLoR/}}
The demo runs in the browser without an installation and saves nothing.
A second mode of the demo opens the schema editor and downloads the project file on save.
The repository also contains a sample screening project and a sample extraction project, a user guide, and the requirement files described in \autoref{sec:requirements}.
